\documentclass[11pt]{article}
\usepackage[T1]{fontenc}
\usepackage{lmodern}
\usepackage{amsmath,amssymb,amsthm,mathtools}
\usepackage{graphicx}
\usepackage[margin=1in]{geometry}
\usepackage{microtype}
\usepackage{xcolor}
\usepackage[colorlinks=true,linkcolor=blue!55!black,citecolor=blue!55!black,urlcolor=blue!55!black]{hyperref}

\newtheorem{theorem}{Theorem}
\newtheorem{lemma}[theorem]{Lemma}
\newtheorem{corollary}[theorem]{Corollary}
\newcommand{\E}{\mathbb E}
\newcommand{\R}{\mathbb R}
\newcommand{\Var}{\operatorname{Var}}
\newcommand{\Cov}{\operatorname{Cov}}
\newcommand{\tr}{\operatorname{tr}}

\title{The General Variance Conjecture from the Isotropic Thin-Shell Theorem}
\author{Literature-implied answer note for arXiv:1209.4270\\
\small Agent: \texttt{agent\_lane\_07} \quad Model: GPT5.6}
\date{22 July 2026}

\begin{document}
\maketitle

\begin{center}
\fbox{\parbox{0.91\textwidth}{
\textbf{Status.} Literature-implied answer (full; agent-identified implication).
Conjectures 1.4 and 1.5 of Alonso-Gutierrez--Bastero are true in all
dimensions. The later supporting paper proves the isotropic theorem but does
not state this anisotropic consequence.
}}
\end{center}

\begin{abstract}
Alonso-Gutierrez and Bastero conjectured that every centered log-concave
random vector $X$ satisfies
\[
  \Var(|X|^2)\leq C\|\Cov(X)\|_{\mathrm{op}}\,\E|X|^2.
\]
Klartag and Lehec recently proved the isotropic variance conjecture in every
dimension. We show by Abel summation over the eigenvalues of $\Cov(X)$ that
the isotropic theorem implies the displayed general inequality with exactly
the same best dimension-free constant. Hence the source's general variance
and general thin-shell conjectures are completely settled. The implication is
elementary but is not stated in the supporting paper found in the bounded
literature search.
\end{abstract}

\section{Source questions}

The source is D. Alonso-Gutierrez and J. Bastero, \emph{The variance
conjecture on some polytopes}, arXiv:1209.4270 \cite{AB}. On PDF page 2,
Conjecture 1.4 asks whether there is a universal $C$ such that every centered
log-concave random vector $X$ satisfies
\begin{equation}\label{eq:general}
  \Var(|X|^2)\leq C\lambda_X^2\E|X|^2,
  \qquad \lambda_X^2=\|\Cov(X)\|_{\mathrm{op}}.
\end{equation}
Conjecture 1.5 asks equivalently whether
\begin{equation}\label{eq:general-shell}
 \sigma_X:=\left[\E\left(|X|-\sqrt{\E|X|^2}\right)^2\right]^{1/2}
 \leq C\lambda_X.
\end{equation}
The authors stress that these statements are not invariant under arbitrary
linear maps and were not known to follow from the isotropic conjecture.

\begin{figure}[ht]
\centering
\includegraphics[width=0.96\textwidth]{figures/open_problem_crop.png}
\caption{Conjectures 1.4 and 1.5 in arXiv:1209.4270, PDF page 2.}
\end{figure}

\section{Later theorem used}

Klartag and Lehec, \emph{Thin-shell bounds via parallel coupling},
arXiv:2507.15495v2 \cite{KL}, prove the following as Theorem 1.1 on PDF
page 2: for every isotropic log-concave random vector $Z$ in $\R^k$,
\begin{equation}\label{eq:isotropic}
  \Var(|Z|^2)=\E(|Z|^2-k)^2\leq C_0 k,
\end{equation}
where $C_0$ is universal. Their statement ranges over every dimension $k$.

\begin{figure}[ht]
\centering
\includegraphics[width=0.96\textwidth]{figures/supporting_answer_crop.png}
\caption{Klartag--Lehec, Theorem 1.1, arXiv:2507.15495v2, PDF page 2.}
\end{figure}

\section{Proof intuition}

Whitening writes the squared norm of an arbitrary centered vector as a
positive quadratic form $\langle MY,Y\rangle$ of an isotropic vector. Order
the eigenvalues of $M$. Abel summation decomposes that quadratic form into a
positive linear combination of the centered energies of nested orthogonal
projections of $Y$. Each projected vector is again isotropic and log-concave,
so the solved isotropic theorem controls each energy in its own dimension.
Minkowski combines these controls without needing any sign information about
mixed fourth moments. A one-line Cauchy--Schwarz estimate for the spectral
counting function then produces exactly $\|M\|_{\mathrm{op}}\tr M$.

The positivity of the spectral increments is essential: it is what permits a
constant-preserving passage from unweighted projection energies to every
positive quadratic form.

\section{Constant-preserving transfer theorem}

For clarity, define the optimal all-dimensional constants
\begin{align*}
C_{\mathrm{iso}}
 &=\sup_{k\geq1}\sup_Z \frac{\Var(|Z|^2)}{k},\\
C_{\mathrm{gen}}
 &=\sup_{n\geq1}\sup_X
 \frac{\Var(|X|^2)}{\|\Cov(X)\|_{\mathrm{op}}\E|X|^2},
\end{align*}
where $Z$ ranges over isotropic log-concave vectors and $X$ ranges over
nonzero centered log-concave vectors. If a law is supported on a proper
linear subspace, all quantities are computed in its linear span.

\begin{theorem}[Exact equivalence]\label{thm:exact}
One has
\[
  C_{\mathrm{gen}}=C_{\mathrm{iso}}.
\]
Consequently, the Klartag--Lehec theorem proves \eqref{eq:general} in every
dimension, with no loss in the universal constant.
\end{theorem}

\begin{proof}
The inequality $C_{\mathrm{iso}}\leq C_{\mathrm{gen}}$ is immediate: for an
isotropic vector in $\R^k$, the denominator in the definition of
$C_{\mathrm{gen}}$ equals $1\cdot k$.

We prove the reverse inequality. Fix a centered log-concave $X$ and reduce to
the linear span of its support. Put $M=\Cov(X)$, which is then positive
definite, and set $Y=M^{-1/2}X$. Thus $Y$ is isotropic and log-concave, and
\[
 |X|^2=\langle MY,Y\rangle,\qquad
 \E|X|^2=\tr M.
\]
Choose an orthonormal eigenbasis of $M$ and write
\[
 \mu_1\geq\mu_2\geq\cdots\geq\mu_n>0,
 \qquad \mu_{n+1}=0.
\]
Let $P_k$ be the orthogonal projection onto the span of the first $k$
eigenvectors. The linear image $P_kY$ is log-concave, and its covariance on
that $k$-dimensional subspace is the identity. Therefore $P_kY$ is isotropic
and, by the definition of $C_{\mathrm{iso}}$,
\begin{equation}\label{eq:partial}
 \left\| |P_kY|^2-k\right\|_{L^2}
 \leq \sqrt{C_{\mathrm{iso}}k}.
\end{equation}

If $Y_i$ are the coordinates of $Y$ in the eigenbasis, Abel summation gives
the exact identity
\begin{align}\label{eq:abel}
 \langle MY,Y\rangle-\tr M
 &=\sum_{i=1}^n\mu_i(Y_i^2-1)\\
 &=\sum_{k=1}^n(\mu_k-\mu_{k+1})
   \sum_{i=1}^k(Y_i^2-1)\notag\\
 &=\sum_{k=1}^n(\mu_k-\mu_{k+1})(|P_kY|^2-k).\notag
\end{align}
All coefficients in the final sum are nonnegative. Minkowski's inequality and
\eqref{eq:partial} yield
\begin{equation}\label{eq:minkowski}
 \sqrt{\Var(\langle MY,Y\rangle)}
 \leq \sqrt{C_{\mathrm{iso}}}
 \sum_{k=1}^n(\mu_k-\mu_{k+1})\sqrt{k}.
\end{equation}

It remains to estimate a deterministic spectral sum. Let
\[
 N(t)=\#\{i:\mu_i\geq t\},\qquad 0\leq t\leq\mu_1.
\]
The layer-cake formula and Cauchy--Schwarz give
\begin{align}\label{eq:spectral}
 \sum_{k=1}^n(\mu_k-\mu_{k+1})\sqrt{k}
 &=\int_0^{\mu_1}\sqrt{N(t)}\,dt\\
 &\leq \left(\mu_1\int_0^{\mu_1}N(t)\,dt\right)^{1/2}
 =\left(\mu_1\sum_{i=1}^n\mu_i\right)^{1/2}.
\end{align}
Combining \eqref{eq:minkowski} with (7)--(8) and squaring proves
\[
 \Var(|X|^2)
 \leq C_{\mathrm{iso}}\,\mu_1\sum_{i=1}^n\mu_i
 =C_{\mathrm{iso}}\|M\|_{\mathrm{op}}\tr M
 =C_{\mathrm{iso}}\lambda_X^2\E|X|^2.
\]
Hence $C_{\mathrm{gen}}\leq C_{\mathrm{iso}}$, completing the proof.
\end{proof}

\begin{corollary}[Source Conjectures 1.4 and 1.5]\label{cor:source}
There is a universal constant $C$ such that every centered log-concave random
vector satisfies both \eqref{eq:general} and \eqref{eq:general-shell}.
\end{corollary}

\begin{proof}
The variance inequality follows from Theorem \ref{thm:exact} and
Klartag--Lehec's bound \eqref{eq:isotropic}. For the thin-shell form, put
$m=\E|X|^2$. Pointwise,
\[
 (|X|-\sqrt m)^2
 \leq \frac{(|X|^2-m)^2}{m},
\]
because $|X|+\sqrt m\geq\sqrt m$. Consequently,
\[
 \sigma_X^2\leq\frac{\Var(|X|^2)}{\E|X|^2}
 \leq C\lambda_X^2.
\]
Taking square roots proves Conjecture 1.5 after renaming the universal
constant.
\end{proof}

\section{Verification and scope}

The proof uses only four points: whitening; closure of log-concavity under
linear images; the Klartag--Lehec theorem in each projected dimension; and the
deterministic Abel identity and spectral estimate displayed above. The accompanying
seeded script checks the Abel identity and spectral inequality on 20,000 random
nonnegative spectra. This computation is an algebra smoke test, not evidence
for the Klartag--Lehec theorem.

The conclusion settles the source's Conjectures 1.4 and 1.5. It does
\emph{not} prove the stronger Kannan--Lovasz--Simonovits conjecture (source
Conjecture 1.3), which concerns every locally Lipschitz test function rather
than only $g(x)=|x|^2$.

\section{Literature and novelty status}

This is classified as a literature-implied answer because the only deep input
is the later theorem \cite{KL}. The supporting authors do not state that they
answer arXiv:1209.4270, and the constant-preserving spectral transfer is the
agent-identified implication.

A bounded search on 22 July 2026 covered the run indexes; the full TeX of
\cite{KL}; exact arXiv/web searches for ``general variance conjecture'', the
displayed anisotropic inequality, and quadratic-form versions; and the later
Alonso-Gutierrez--Bastero paper arXiv:1610.04023 \cite{ABlater}. The latter
still says that equivalence between the isotropic and general conjectures was
unclear and proves special families. No later paper explicitly stating
Theorem \ref{thm:exact} or the full anisotropic conclusion was found. Because
the transfer lemma is short, it may nevertheless be folklore; novelty
confidence is moderate.

\section{Human review recommendation}

Treat this as a full literature-implied answer. Reviewers should verify the
constant-preserving Abel decomposition, the use of the isotropic theorem on
all orthogonal marginals, and the provenance/novelty classification. The
mathematical scope is all centered log-concave laws and includes degenerate
laws after restriction to their linear span.

\begin{thebibliography}{9}
\bibitem{AB}
D. Alonso-Gutierrez and J. Bastero,
\emph{The variance conjecture on some polytopes},
arXiv:1209.4270 (2012; PDF revision dated 2018).

\bibitem{KL}
B. Klartag and J. Lehec,
\emph{Thin-shell bounds via parallel coupling},
arXiv:2507.15495v2 (2026), Theorem 1.1.

\bibitem{ABlater}
D. Alonso-Gutierrez and J. Bastero,
\emph{The variance conjecture on hyperplane projections of the
$\ell_p^n$ balls},
arXiv:1610.04023; Rev. Mat. Iberoam. 34 (2018), 879--904.
\end{thebibliography}

\end{document}
